\part{检索增强生成与智能体}

\chapter{检索增强生成（RAG）}\label{sec:rag}

\begin{motivationbox}{参数记忆无法及时更新，也难以给出可核验来源}
预训练知识被压缩在参数中，更新一次通常需要重新训练或微调；面对企业私有数据、刚发生的事件和要求精确引用的任务，单靠参数记忆既可能过时，也难以追溯证据。检索增强生成（Retrieval-Augmented Generation，RAG）在生成前动态检索外部资料，把“记住所有事实”改成“需要时找到事实并放入上下文”。
\end{motivationbox}

RAG 用外部检索补充模型上下文，常见流程是：文档切块、计算向量表示（embedding）、写入向量库、查询召回、重新排序、拼接上下文、生成答案。

\section{为什么需要切块、召回与重排}

整篇文档通常长于向量模型或生成模型的有效输入，也很难用一个向量表达所有局部主题，因此首先需要\textbf{文档切块（chunking）}。块太大时，向量语义会被多个主题稀释，命中后还会带入大量噪声；块太小时，句子和表格上下文会被切断。重叠窗口、按标题分段和保留父文档元数据，都是为了缓解边界信息丢失。

向量召回的目标是从海量文档中低成本取得较高\textbf{召回率（recall）}，即尽量不要漏掉相关材料，因此常使用近似最近邻，并容忍候选中混入噪声。仅靠向量相似度又难以精细判断复杂查询与段落的逐词对应，于是增加\textbf{重排器（reranker）}，用计算更贵的交叉编码或模型评分，在小候选集上提高\textbf{精确率（precision）}，也就是减少无关材料。两阶段设计本质上是“便宜地广搜，再昂贵地精排”。

最终只把有限数量的高价值片段放入上下文，是因为上下文窗口有限，提示词预填充（prefill）的计算成本也会随输入变长，而且更多材料不一定更好。噪声证据会分散注意力，互相矛盾的片段还可能诱导错误回答。因此 RAG 的关键权衡是召回率、精确率、上下文预算与延迟，而不是简单增加返回候选数 Top-$k$。

稠密向量检索擅长语义近似，BM25 等词法检索更擅长精确名称、编号和罕见词；\textbf{混合检索（hybrid search）}把二者候选或分数融合。元数据过滤则在检索前限制租户、权限、语言、产品版本和生效日期。它不是“锦上添花”：当旧版制度与新版制度语义高度相似时，单靠 embedding 很难可靠判断哪个版本有效。

长上下文还存在 \textbf{Lost in the Middle} 风险：相关证据虽然被放入窗口，却位于大量文本中间而没有被模型稳定利用。改善方式包括重排、去重、压缩、把高置信证据放在更显著的位置，以及要求引用与证据核验；简单扩大 Top-$k$ 往往同时增加预填充成本与干扰。

RAG 不能从根本上消除幻觉，因为模型仍可能误读、忽略或编造上下文。要提高可信度，还需要要求答案引用检索片段、检查引用是否真正支持结论，并在证据不足时允许拒答。RAG 解决知识访问问题，不自动解决推理正确性和生成忠实性。

\begin{examplebox}{查询一条刚更新的公司制度}
用户询问：“2026 年正式员工有多少天年假？”知识库中有三个候选块：
\begin{itemize}
  \item A：2024 年旧制度，年假 10 天；
  \item B：2026 年新制度，年假调整为 15 天；
  \item C：2026 年差旅报销标准。
\end{itemize}
向量召回可能因“年假”相似度返回 $[A,B,C]$。重排器同时阅读问题与段落后，应把包含正确年份和主题的 B 排在最前。生成上下文只放 B，并保留文档日期与来源，模型才能回答“15 天”并引用新制度。

若系统只按语义相似度、不检查版本元数据，旧制度 A 也很可能排名很高；若 B 根本没有被召回，生成模型再强也无法从上下文恢复缺失事实。这个例子分别对应检索的召回率、重排的精确率和版本过滤问题。
\end{examplebox}

\begin{takeawaybox}
RAG 面试不要只说“接向量数据库”。要能解释文档切块、向量质量、召回评估、重新排序、上下文窗口和引用可信度。
\end{takeawaybox}

\begin{interviewbox}{为什么高召回之后仍需要 reranker？}
\textbf{20 秒回答：}第一阶段检索要在大库中便宜地减少漏召，通常允许较多噪声；reranker 在小候选集上做更细的查询--文档联合判断，提高精确率。二者分别优化候选覆盖和最终上下文质量。

\textbf{深入解释：}双编码器可以预先计算文档向量，检索成本低，但查询与文档只在最终向量上交互；交叉编码器能逐词联合建模，判断更精细，却不适合扫描全库。若正确片段没有进入第一阶段候选，reranker 无法恢复它；若把所有高召回候选直接塞给模型，噪声和上下文成本又会上升。

\textbf{面试官继续追问：}
1. chunk 太大和太小分别损失什么？
2. metadata filtering 应发生在召回前还是召回后？
3. hybrid search 怎样合并不同分数尺度？
4. 如何分别评估 retrieval 与 generation？
\end{interviewbox}

\begin{interviewbox}{上下文越多为什么不一定越好？}
\textbf{20 秒回答：}更多上下文增加预填充成本，也会引入重复、过时或冲突证据；模型还可能忽略长上下文中间的信息。目标应是证据覆盖率与信噪比，而不是最大化输入长度。

\textbf{深入解释：}版本冲突需要以来源、生效时间和权限等元数据解决，不能让生成模型自行“投票”。系统还应在文档更新时保存版本标识，在答案中引用实际使用的片段，并对无充分证据的查询允许拒答。

\textbf{面试官继续追问：}
1. Lost in the Middle 是检索失败还是利用失败？
2. 文档版本冲突如何构造回归测试？
3. 为什么引用存在不等于引用支持结论？
\end{interviewbox}

\chapter{智能体（Agent）}\label{sec:agents}

\begin{motivationbox}{一次模型调用无法完成需要外部反馈的多步任务}
普通 LLM 前向只能根据当前上下文生成文本，无法知道数据库查询是否成功、代码执行结果是什么，也不能在看到新结果后自动调整计划。智能体（Agent）把模型置于“观察--行动”（observation--action）闭环中：模型选择工具，环境返回真实观察，模型再据此继续决策，从而处理需要外部状态和多步操作的任务。
\end{motivationbox}

智能体通常包含模型、工具、状态/记忆、规划和观察--行动循环。一个简化循环是：
\[
\text{observe}\rightarrow\text{plan}\rightarrow\text{act(tool)}\rightarrow\text{observe}.
\]
\section{为什么需要参数规范、状态与安全边界}

\textbf{工具参数规范（schema）}把自然语言意图约束为结构化参数，解决模型输出难以被程序可靠解析的问题。参数规范越清楚，参数类型、必填项和错误反馈越明确，模型越容易选择正确工具；但它不能代替服务端校验，模型仍可能提供不存在的标识符（ID）或越权参数。

\textbf{状态与记忆} 用于保存跨步骤事实、计划进度和工具结果，解决有限上下文下的信息持续性问题。原样保留全部历史会快速耗尽窗口并增加噪声，因此需要摘要、结构化状态或按需检索。压缩又可能丢失关键细节，所以重要状态应有明确数据结构，而不只依赖自然语言摘要。

\textbf{权限边界与确认}是因为智能体的错误不再只是一段错误文本，还可能变成真实写操作。读取、写入、付款、删除等工具应使用最小权限，对高影响或不可逆动作增加参数校验、幂等设计和人工确认。错误恢复则要区分可重试的临时失败与逻辑错误，避免模型无限循环调用。

\textbf{可观测性}记录每一步输入、工具选择、参数、响应和延迟，用于定位长链路中的首个错误。智能体的主要风险正是误差累积：一步错误观察或工具调用会污染后续计划；如果只看最终回答，通常无法判断问题来自模型、工具、数据还是权限层。

\textbf{幂等性（idempotency）}指同一操作重复执行时，最终效果与执行一次一致。网络超时后客户端未必知道服务端是否已经完成写入；若没有幂等键就直接重试，可能重复付款或创建两场会议。对无法天然幂等的动作，应使用唯一请求标识、先查询后重试或事务状态机。

\textbf{MCP（Model Context Protocol）}在面试层面可以理解为标准化的上下文、工具和资源接口。模型应用负责推理与编排，外部能力提供方按统一协议暴露可发现的工具或资源，从而降低每个应用重复编写私有适配器的成本。MCP 解决接口与集成问题，不自动解决权限、正确性、工具选择和业务安全。

RAG 和智能体的关系也应区分：RAG 通常是一次生成前的知识检索流程，智能体则允许模型根据观察反复决定下一步动作。检索可以成为智能体的一个工具，但使用了向量数据库并不意味着系统已经具备规划和反馈闭环。

\begin{examplebox}{安排一次不能冲突的会议}
用户要求“给项目组安排明天下午 30 分钟会议”。一个可靠智能体的执行链可以是：
\begin{enumerate}
  \item 从用户请求中提取参会人、日期范围和时长；
  \item 调用日历读取工具，观察每个人的忙闲状态；
  \item 发现 14:00 有人冲突后，重新规划并选择共同空闲的 15:30；
  \item 调用会议室查询工具，得到可用房间 R2；
  \item 在真正创建事件前向用户确认“15:30--16:00，R2”；
  \item 用户确认后调用写入工具，并读取返回的事件编号验证成功。
\end{enumerate}
普通一次函数调用通常只能完成第 6 步的参数提交；智能体的关键是前面的多次观察、根据冲突修改计划，以及对写操作设置确认边界。若创建接口超时，还应先用事件编号或幂等键查询结果，而不是盲目重试并创建重复会议。
\end{examplebox}

\begin{interviewbox}{Agent 和 workflow 的本质区别是什么？}
\textbf{20 秒回答：}workflow 的步骤和分支主要由程序预先定义；Agent 把下一步动作选择部分交给模型，并根据环境观察动态循环。实际系统常采用受约束的混合模式，而不是让模型无限自由规划。

\textbf{深入解释：}固定流程更可预测、易测试；开放式 Agent 能处理未预先枚举的状态，但增加循环、权限和错误累积风险。规划也不一定要输出可见的长文本，可以体现为结构化状态更新、策略选择或下一工具调用。

\textbf{面试官继续追问：}
1. Agent 为什么需要显式状态？
2. 哪些任务更适合固定 workflow？
3. 如何设置最大步数和停止条件？
4. tool error 应该重试、换工具还是交给用户？
\end{interviewbox}

\begin{interviewbox}{为什么 tool call 需要 schema、确认与幂等性？}
\textbf{20 秒回答：}schema 让参数可解析和校验；确认把高影响决策留给用户；幂等性避免超时重试产生重复副作用。三者分别控制输入格式、授权边界和执行语义。

\textbf{深入解释：}schema 不能代替服务端鉴权，确认也不能修复错误参数。可靠执行还需要最小权限、超时分类、幂等键、审计日志和补偿机制。MCP 可以统一能力接口，但这些业务约束仍由应用与服务端落实。

\textbf{面试官继续追问：}
1. 为什么 destructive action 必须二次确认？
2. Agent loop 为什么可能无限循环？
3. MCP 与具体工具实现之间如何解耦？
4. 什么情况下重试反而危险？
\end{interviewbox}

\section*{面试高频问题}
\addcontentsline{toc}{section}{面试高频问题：RAG 与 Agents}

\textbf{Q: RAG 的召回率和精确率如何影响答案？}

召回率低会漏掉关键事实；精确率低会把噪声塞进上下文，导致模型被干扰或引用错误材料。

\textbf{Q: 智能体和普通函数调用有什么不同？}

智能体会根据观察迭代规划和调用工具；普通函数调用通常是一次性参数填充，不一定有多步状态和反馈循环。
